% Appendix: Thermodynamic Derivations
% Converted on 2026-10-08 from docs/lswt/04-appendices/thermodynamic-derivations.md (status: draft, source-section: "Thermodynamics in Linear Spin Wave Theory; Number and Spin moment from Correlation Matrix (source pp. 13–16); restructured TeX appendix 'Thermodynamic derivations'").
% Review status: draft; awaiting the user's physics and mathematics review.

\hypertarget{sec-lswt-thermodynamic-derivations}{%
\section{Appendix: Thermodynamic Derivations}\label{sec-lswt-thermodynamic-derivations}}

The thermodynamic quantities of a harmonic magnon gas and the equal-time correlation matrix of the Holstein--Primakoff bosons both follow from the Gibbs state of the diagonal magnon Hamiltonian. Thermodynamics (Section~\ref{sec-lswt-thermodynamics}) and Magnon Observables (Section~\ref{sec-lswt-magnon-observables}) state the results with short derivations; this appendix gives the intermediate steps: the factorization of the partition function, the entropy and heat-capacity identities, the small-energy expansions behind the zero-mode table, and the correlation-matrix route to the boson number. We assume, as in those documents, that \(\mathsf H_{\mathbf k}\) is positive definite at every momentum of the magnetic Brillouin zone (MBZ), so that every magnon energy \(\varepsilon_{n\mathbf k}\) is positive, except where the zero-mode limit is discussed.

\subsection{Factorization of the Partition Function}

The diagonal Hamiltonian \(\hat H_{\mathrm{LSWT}}=E_{\mathrm{GS}}+\sum_{\mathbf k,n}\varepsilon_{n\mathbf k}\hat n_{\mathbf kn}\) of \lswtnoteref{LSWT paraunitary diagonalization}{Paraunitary Diagonalization} is diagonal in the Fock basis \(|\{m_{\mathbf kn}\}\rangle\) of magnon occupation numbers \(m_{\mathbf kn}=0,1,2,\ldots\), with eigenvalues \(E_{\mathrm{GS}}+\sum_{\mathbf k,n}\varepsilon_{n\mathbf k}m_{\mathbf kn}\). The trace is a sum over all occupation configurations, and the exponential of a sum is a product, so

\[
Z=\exp(-\beta E_{\mathrm{GS}})\prod_{\mathbf k,n}\sum_{m=0}^{\infty}\exp(-\beta\varepsilon_{n\mathbf k}m)
=\exp(-\beta E_{\mathrm{GS}})\prod_{\mathbf k,n}\frac{1}{1-\exp(-\beta\varepsilon_{n\mathbf k})}.
\]

Each geometric series converges because \(\varepsilon_{n\mathbf k}>0\). The Gibbs state is the product of the single-mode thermal states \(\hat\rho_{\mathbf kn}=(1-\exp(-\beta\varepsilon_{n\mathbf k}))\exp(-\beta\varepsilon_{n\mathbf k}\hat n_{\mathbf kn})\). The occupation of one mode follows from the derivative of the corresponding factor, \(\langle\hat n_{\mathbf kn}\rangle=-\beta^{-1}\partial\ln Z/\partial\varepsilon_{n\mathbf k}=[\exp(\beta\varepsilon_{n\mathbf k})-1]^{-1}=n_{\mathrm B}(\varepsilon_{n\mathbf k})\).

\subsection{Internal Energy and Free Energy}

The internal energy is \(U=\operatorname{Tr}(\hat\rho\hat H_{\mathrm{LSWT}})=-\partial\ln Z/\partial\beta\). Differentiating the logarithm of each factor gives

\[
-\frac{\partial}{\partial\beta}\ln\frac{1}{1-\exp(-\beta\varepsilon)}
=\frac{\varepsilon\exp(-\beta\varepsilon)}{1-\exp(-\beta\varepsilon)}
=\varepsilon\,n_{\mathrm B}(\varepsilon),
\]

and the prefactor contributes \(E_{\mathrm{GS}}\), which yields \(U=E_{\mathrm{GS}}+\sum_{\mathbf k,n}\varepsilon_{n\mathbf k}n_{\mathrm B}(\varepsilon_{n\mathbf k})\). The free energy \(F=-\beta^{-1}\ln Z\) follows by taking the logarithm of the product, \(F=E_{\mathrm{GS}}+\beta^{-1}\sum_{\mathbf k,n}\ln(1-\exp(-\beta\varepsilon_{n\mathbf k}))\).

\subsection{Entropy}

The von Neumann entropy of the Gibbs state uses \(\ln\hat\rho=-\beta\hat H_{\mathrm{LSWT}}-\ln Z\):

\[
\mathcal S=-k_{\mathrm B}\operatorname{Tr}(\hat\rho\ln\hat\rho)
=k_{\mathrm B}\beta\operatorname{Tr}(\hat\rho\hat H_{\mathrm{LSWT}})+k_{\mathrm B}\ln Z
=\frac{U-F}{T}.
\]

The thermodynamic entropy \(-\partial F/\partial T\) gives the same result. With \(F=-k_{\mathrm B}T\ln Z\) and \(\partial\beta/\partial T=-1/(k_{\mathrm B}T^2)\),

\[
-\frac{\partial F}{\partial T}
=k_{\mathrm B}\ln Z+k_{\mathrm B}T\,\frac{\partial\ln Z}{\partial\beta}\frac{\partial\beta}{\partial T}
=k_{\mathrm B}\ln Z+\frac{U}{T}
=\frac{U-F}{T}.
\]

For the Gibbs state the von Neumann entropy therefore reduces to the thermodynamic entropy. Inserting the expressions for \(U\) and \(F\), the constant \(E_{\mathrm{GS}}\) cancels in \(U-F\), and

\[
\mathcal S=k_{\mathrm B}\sum_{\mathbf k,n}\Big[\beta\varepsilon_{n\mathbf k}n_{\mathbf kn}-\ln\big(1-\exp(-\beta\varepsilon_{n\mathbf k})\big)\Big],
\qquad
n_{\mathbf kn}=n_{\mathrm B}(\varepsilon_{n\mathbf k}).
\]

The definition \(n=[\exp(\beta\varepsilon)-1]^{-1}\) gives \(\exp(\beta\varepsilon)=(1+n)/n\) and \(1-\exp(-\beta\varepsilon)=1/(1+n)\), so

\[
\beta\varepsilon=\ln\frac{1+n}{n},
\qquad
\ln\big(1-\exp(-\beta\varepsilon)\big)=-\ln(1+n).
\]

Substituting these relations, the summand becomes \(n\ln(1+n)-n\ln n+\ln(1+n)=(1+n)\ln(1+n)-n\ln n\), which is the occupation form of the entropy in Thermodynamics (Section~\ref{sec-lswt-thermodynamics}).

\subsection{Heat Capacity}

The temperature derivative of the occupation is

\[
\frac{\partial n_{\mathrm B}(\varepsilon)}{\partial T}
=\frac{\varepsilon}{k_{\mathrm B}T^2}\,\frac{\exp(\beta\varepsilon)}{\big(\exp(\beta\varepsilon)-1\big)^2}
=\frac{\varepsilon}{k_{\mathrm B}T^2}\,\frac{1}{4\sinh^2(\beta\varepsilon/2)},
\]

where the second form uses \(\exp(x)/(\exp(x)-1)^2=1/(\exp(x/2)-\exp(-x/2))^2\). Since only the occupations depend on temperature, \(C=\partial U/\partial T=\sum_{\mathbf k,n}\varepsilon_{n\mathbf k}\,\partial n_{\mathbf kn}/\partial T\), which gives the expression \(k_{\mathrm B}\sum(\beta\varepsilon)^2/[4\sinh^2(\beta\varepsilon/2)]\).

The relation \(C=T\,\partial\mathcal S/\partial T\) follows from the occupation form of the entropy. The derivative of the summand with respect to the occupation is \(k_{\mathrm B}\ln[(1+n)/n]=k_{\mathrm B}\beta\varepsilon\), so

\[
T\frac{\partial\mathcal S}{\partial T}
=T\sum_{\mathbf k,n}k_{\mathrm B}\beta\varepsilon_{n\mathbf k}\,\frac{\partial n_{\mathbf kn}}{\partial T}
=\sum_{\mathbf k,n}\varepsilon_{n\mathbf k}\frac{\partial n_{\mathbf kn}}{\partial T}
=\frac{\partial U}{\partial T}.
\]

\subsection{Small-Energy Expansions Behind the Zero-Mode Limits}

With \(x=\beta\varepsilon\to0^+\), the summands of the thermodynamic quantities have the expansions

\[
\begin{aligned}
\varepsilon\,n_{\mathrm B}(\varepsilon)&=k_{\mathrm B}T\left(1-\frac x2+\frac{x^2}{12}+O(x^4)\right),\\
k_{\mathrm B}T\ln\big(1-\exp(-x)\big)&=k_{\mathrm B}T\left(\ln x-\frac x2+\frac{x^2}{24}+O(x^4)\right),\\
k_{\mathrm B}\Big[x\,n_{\mathrm B}-\ln\big(1-\exp(-x)\big)\Big]&=k_{\mathrm B}\left(1-\ln x+\frac{x^2}{24}+O(x^4)\right),\\
k_{\mathrm B}\frac{x^2}{4\sinh^2(x/2)}&=k_{\mathrm B}\left(1-\frac{x^2}{12}+O(x^4)\right).
\end{aligned}
\]

They follow from \(x/(\exp(x)-1)=1-x/2+x^2/12+O(x^4)\) and \(\ln(1-\exp(-x))=\ln x-x/2+x^2/24+O(x^4)\). The energy and heat-capacity summands stay finite, while the free-energy and entropy summands diverge as \(\ln x=\ln(\varepsilon/k_{\mathrm B}T)\). Near an isolated zero \(\varepsilon_{n\mathbf k}\propto|\mathbf k-\mathbf k_0|^a\) with \(a>0\), the divergent part of each summand is proportional to \(a\ln|\mathbf k-\mathbf k_0|\), and in two dimensions \(\int_{|\mathbf q|<q_0}d^2\mathbf q\,\ln|\mathbf q|=\pi q_0^2(\ln q_0-\frac12)\) is finite. The momentum integrals of \(F\) and \(\mathcal S\) therefore remain finite in the thermodynamic limit, as stated in Thermodynamics (Section~\ref{sec-lswt-thermodynamics}). Finiteness of these integrals does not imply that the magnon-gas description is controlled; that condition involves the boson number, treated below.

\subsection{Correlation Matrix of the Holstein--Primakoff Bosons}

The equal-time correlation matrix \(\langle\hat\Psi_{\mathbf k}\hat\Psi_{\mathbf k}^\dagger\rangle\) contains every average of two Holstein--Primakoff operators at momentum \(\mathbf k\). Its blocks are

\[
\big\langle\hat\Psi_{\mathbf k}\hat\Psi_{\mathbf k}^\dagger\big\rangle
=\begin{pmatrix}
\langle\hat{\mathbf a}_{\mathbf k}\hat{\mathbf a}_{\mathbf k}^\dagger\rangle & \langle\hat{\mathbf a}_{\mathbf k}\hat{\mathbf a}_{-\mathbf k}\rangle\\
\langle\hat{\mathbf a}_{-\mathbf k}^\dagger\hat{\mathbf a}_{\mathbf k}^\dagger\rangle & \langle\hat{\mathbf a}_{-\mathbf k}^\dagger\hat{\mathbf a}_{-\mathbf k}\rangle
\end{pmatrix},
\]

where each block is the \(N_{\mathrm{sub}}\times N_{\mathrm{sub}}\) matrix of averages over the sublattice indices. The average is linear, and the Bogoliubov transformation \(\hat\Psi_{\mathbf k}=\mathsf T_{\mathbf k}\hat\Phi_{\mathbf k}\) has c-number coefficients, so

\[
\big\langle\hat\Psi_{\mathbf k}\hat\Psi_{\mathbf k}^\dagger\big\rangle
=\mathsf T_{\mathbf k}\big\langle\hat\Phi_{\mathbf k}\hat\Phi_{\mathbf k}^\dagger\big\rangle\mathsf T_{\mathbf k}^\dagger.
\]

The Gibbs state is diagonal in the magnon occupation numbers, so every average of two magnon operators that changes some occupation vanishes. This removes \(\langle\hat b\hat b\rangle\), \(\langle\hat b^\dagger\hat b^\dagger\rangle\), and the products of different modes. The remaining averages are \(\langle\hat b_{\mathbf kn}\hat b_{\mathbf kn}^\dagger\rangle=1+n_{n\mathbf k}\) and \(\langle\hat b_{-\mathbf kn}^\dagger\hat b_{-\mathbf kn}\rangle=n_{n,-\mathbf k}\), where \(n_{n\mathbf k}=n_{\mathrm B}(\varepsilon_{n\mathbf k})\). Hence

\[
\big\langle\hat\Phi_{\mathbf k}\hat\Phi_{\mathbf k}^\dagger\big\rangle
=\operatorname{diag}\big(1+n_{1\mathbf k},\ldots,1+n_{N_{\mathrm{sub}}\mathbf k},\;n_{1,-\mathbf k},\ldots,n_{N_{\mathrm{sub}},-\mathbf k}\big),
\]

which is the matrix \(\mathsf N_{\mathbf k}(0)\) of \lswtnoteref{LSWT spin correlations}{Spin Correlations}; each hole entry carries the occupation of its own band at \(-\mathbf k\). In the magnon vacuum all occupations vanish and \(\mathsf N_{\mathbf k}(0)=\operatorname{diag}(\mathsf I_{N_{\mathrm{sub}}},0)\).

\subsection{Boson Number and Sublattice Moment}

The boson number of sublattice \(\mu\) averaged over the \(N_{\mathrm{uc}}\) magnetic unit cells is \(\langle\hat n_\mu\rangle=N_{\mathrm{uc}}^{-1}\sum_i\langle\hat a_{i\mu}^\dagger\hat a_{i\mu}\rangle\). The Fourier phases \(\exp(\pm\mathrm i\mathbf k\cdot\mathbf r_{i\mu})\) of the full-position convention cancel in the product of an operator and its adjoint on the same sublattice, and the sum over cells gives \(\sum_i\hat a_{i\mu}^\dagger\hat a_{i\mu}=\sum_{\mathbf k}\hat a_{\mathbf k\mu}^\dagger\hat a_{\mathbf k\mu}\). The diagonal entry \(N_{\mathrm{sub}}+\mu\) of the correlation matrix at \(\mathbf k\) is \(\langle\hat a_{-\mathbf k\mu}^\dagger\hat a_{-\mathbf k\mu}\rangle\). The lower block row of \(\mathsf T_{\mathbf k}\) is \((\mathsf Q_{\mathbf k}^*,\mathsf P_{-\mathbf k}^*)\), so

\[
\Big[\mathsf T_{\mathbf k}\mathsf N_{\mathbf k}(0)\mathsf T_{\mathbf k}^\dagger\Big]_{N_{\mathrm{sub}}+\mu,\,N_{\mathrm{sub}}+\mu}
=\sum_{n=1}^{N_{\mathrm{sub}}}\Big[\big|(\mathsf Q_{\mathbf k})_{\mu n}\big|^2\big(1+n_{n\mathbf k}\big)+\big|(\mathsf P_{-\mathbf k})_{\mu n}\big|^2n_{n,-\mathbf k}\Big].
\]

Summing over a momentum set that is closed under inversion and relabeling \(-\mathbf k\to\mathbf k\) in the \(\mathsf P\) term gives the boson number of Magnon Observables (Section~\ref{sec-lswt-magnon-observables}),

\[
\langle\hat n_\mu\rangle=\frac{1}{N_{\mathrm{uc}}}\sum_{\mathbf k\in\mathrm{MBZ}}\sum_{n=1}^{N_{\mathrm{sub}}}\Big[\big|(\mathsf P_{\mathbf k})_{\mu n}\big|^2n_{n\mathbf k}+\big|(\mathsf Q_{\mathbf k})_{\mu n}\big|^2\big(1+n_{n\mathbf k}\big)\Big].
\]

The same result follows directly from \(\hat a_{\mathbf k\mu}=\sum_n[(\mathsf P_{\mathbf k})_{\mu n}\hat b_{\mathbf kn}+(\mathsf Q_{-\mathbf k})_{\mu n}\hat b_{-\mathbf kn}^\dagger]\) and the magnon averages above. The longitudinal spin component in the local frame is \(\hat{\widetilde S}_{i\mu}^0=S_\mu-\hat a_{i\mu}^\dagger\hat a_{i\mu}\) to the order of linear spin-wave theory, so the cell-averaged longitudinal moment of sublattice \(\mu\) is \(S_\mu-\langle\hat n_\mu\rangle\).

\subsection{References}

\subsubsection{Internal Documents}

\begin{itemize}
\tightlist
\item
  Thermodynamics (Section~\ref{sec-lswt-thermodynamics}): states the thermodynamic quantities and the zero-mode limits derived here.
\item
  Magnon Observables (Section~\ref{sec-lswt-magnon-observables}): states the correlation matrix, the boson number, and the reduced moment derived here.
\item
  \lswtnoteref{LSWT spin correlations}{Spin Correlations}: defines the magnon correlation matrix \(\mathsf N_{\mathbf q}(t)\), whose equal-time value is derived here.
\item
  \lswtnoteref{LSWT paraunitary diagonalization}{Paraunitary Diagonalization}: supplies the diagonal magnon Hamiltonian and the blocks of \(\mathsf T_{\mathbf k}\).
\item
  \lswtnoteref{LSWT boson Hamiltonian}{Momentum-Space BdG Hamiltonian}: defines the full-position Fourier convention and the Nambu spinor.
\item
  \lswtnoteref{LSWT spin Hamiltonian}{Notation and Conventions}: fixes the thermodynamic and system-size notation.
\end{itemize}

\subsubsection{External Sources}

\begin{itemize}
\tightlist
\item
  None.
\end{itemize}
